\section{Experimental setup}\label{sec:setup}

\subsection{Settings}\label{sec:settings}

We evaluate the methods in two settings. Setting~A holds all \NCases{} evaluation cases of the three
families, and Setting~B the \NSynth{} synthesized cases alone, 100 for each number of ground-truth
equations from 1 to 10. Every result uses the stratified source-disjoint split of
Section~\ref{sec:split} with the same 10 random seeds.

\subsection{Compared methods}\label{sec:compared}

We compare the proposed method with a classical information-retrieval baseline, two variants of the
baseline with stronger text scorers, and two LLMs. Table~\ref{tab:inputs} lists the information each
method uses.

\begin{table}[tbp]
\centering
\caption{Information used by each method. The baseline includes its variants with BM25 and E5.}
\label{tab:inputs}
\footnotesize
\setlength{\tabcolsep}{4pt}
\begin{tabular}{p{0.44\linewidth}ccc}
\toprule
Information & Proposed & Baseline & LLMs \\
\midrule
Query (description, input and output variables) & yes & yes & yes \\
Expression and domain label of each equation & yes & yes & yes \\
Variable symbols of each equation, as a list & yes & yes & no \\
Context text of each equation & yes & yes & no \\
Source of each equation & no & no & no \\
Cases with known ground truth & training, validation & validation & none \\
\bottomrule
\end{tabular}
\end{table}

\paragraph{Classical baseline} The baseline ranks the whole database by the Stage-1 score of
Eq.~\eqref{eq:stage1}, without a learned model and without a candidate set. Its weights decide how
well it ranks, so we compare \NBasePairs{} weight pairs $(w_{\mathrm{t}}, w_{\mathrm{v}})$ from
$\BasePairFirst$ to $\BasePairLast$ and, in each split, evaluate on the test cases the pair with the highest Recall@20 on
the validation cases (\ref{app:weights}). Two variants test whether stronger text evidence alone
would close the gap to the proposed method. The first replaces the TF-IDF cosine with BM25
\citep{BM25}; the second replaces it with the cosine between dense embeddings of the description
and of each equation text from the sentence encoder E5 \citep{E5}. Both rescale the text score to
$[0, 1]$ over the database for every query, choose their weights on the validation cases in the
same way, and train nothing.

\paragraph{LLMs} An LLM receives the whole database and the query of one test case and returns the
identifiers of 20 equations for the desired model, from the most to the least likely. It sees no
solved case (zero-shot). Each equation reaches it only as its identifier, its expression, and its
domain label: the prompt with the database holds \NDbTokens{} tokens, and adding the variable
definitions and the context texts would exceed the input limit of one million tokens. We withhold
the sources, shuffle the order of the equations, and renumber their identifiers, because the
proposed method does not read the sources and the synthesized ground truths prefer equations of one
source, so the sources would give the LLMs alone a clue to the answer. The prompt states the first
three requirements of Section~\ref{sec:formulation} (\ref{app:prompts}). The models are Claude
Opus~5.5 and Claude Haiku~5.5 \citep{ClaudeModels}, a large and a small model of one vendor, both
with the reasoning effort set to high. An LLM sees no training case, so its input for a case does
not depend on the split; we therefore ask it three times per case, independently, and use the mean
of the three answers in every split in which the case is a test case. The splits of the two
settings hold \NTestDistinct{} distinct test cases, so each model answered \NCalls{} calls.

\paragraph{Proposed method} We compared four pairs of first-stage weights, \StagePairs, by the
Recall@20 of the validation cases. Their validation
Recall@20 ranged only from \ValRecAMin{} to \ValRecAMax{} in Setting~A and from \ValRecBMin{} to
\ValRecBMax{} in Setting~B, so we use the pair $(\WText, \WVar)$ in every split, whose candidate set
admits every ground-truth equation with the smallest $N_1$. In each split we set $N_1$ to the
smallest number of candidates that holds every ground-truth equation of the training and validation
cases, rounded up to the next thousand; this gives $N_1 = \NOne$ in all 10 splits, whereas the
other pairs need \NOneOtherMin{} to \NOneOtherMax. Between random and hard negatives, each split
takes the one with the higher validation Recall@20, and hard negatives won in all \NSplitsAll{}
splits of the two settings. The other hyperparameters are a reference set of $m = \RefSize$
equations, the margin $\gamma = \Margin$, $n_{\mathrm{neg}} = \NNeg$ incorrect candidates, and
\Epochs{} epochs of AdamW with learning rate $\LearnRate$; we have not compared them with other
values.

\subsection{Judgment and metrics}\label{sec:metrics}

Every method returns identifiers of database equations: the retrieval methods rank the database,
and the LLMs list identifiers. A returned equation is correct when its identifier belongs to the
ground truth $\mathcal{R}_q$; an identifier absent from the database counts as incorrect and keeps
its rank, and a repeated identifier counts only at its first occurrence. For each test case we
compute the following quantities and average them over the test cases of a split;
$\Top_n(q)$ denotes the top $n$ equations that a method returns.
\begin{itemize}
  \item Recall@20 $= |\mathcal{R}_q \cap \Top_{20}(q)| / K$, the primary metric: the fraction of
    the ground-truth equations among the 20 equations presented.
  \item Mean average precision (MAP): the mean over cases of the average precision
    $\mathrm{AP}(q) = \frac{1}{K} \sum_{e \in \mathcal{R}_q} P(e)$, where $P(e)$ is the fraction of
    ground-truth equations among the equations ranked at or above $e$, and $P(e) = 0$ when $e$ is
    not in the ranking. MAP shows how closely the ground-truth equations crowd at the top. Because
    an LLM returns only 20 equations, the comparison with LLMs uses MAP@20, in which $P(e) = 0$ for
    every ground-truth equation outside the top 20.
  \item R-Precision $= |\mathcal{R}_q \cap \Top_K(q)| / K$: the fraction of the ground-truth
    equations within the top $K$ \citep{Manning}.
\end{itemize}
We run no statistical tests. To compare two methods, we report the mean of their difference over
the 10 splits and the number of splits in which one beats the other.

\subsection{Implementation}\label{sec:impl}

Every experiment runs as a deterministic script: the random seed fixes the split, the model
initialization, and the choice of incorrect candidates. The proposed method calls no language model
at inference. One run of training and evaluation for a split of Setting~A took \TimeRunA{} seconds on
average, with six runs in parallel on the CPU of a laptop computer (\ref{app:n1}).
